% GENERATED FILE. Edit canonical lessons, then run npm run generate:books.
% canonical-source-hash: fnv1a64:f5130c5743833c5f
% canonical-lessons: RU-C409-polomka, RU-C409-evakuator, RU-C409-buksirovat, RU-C409-domkrat, RU-C409-avtoservis

\chapter{Breaking down}
\label{ch:ru-breaking-down}
\hlchaptermodality{409}

\begin{chapteropening}
\textbf{By the end of this chapter:} \emph{I can get help after a breakdown: the breakdown, the tow truck, towing, the jack and the car repair shop.}

Complete the last lesson of chapter 409: I can get help after a breakdown: the breakdown, the tow truck, towing, the jack and the car repair shop.
\end{chapteropening}

\section[\texorpdfstring{\ru{поломка} (polómka)}{поломка (polómka)}]{\ru{поломка} (polómka) — a breakdown, a fault}

\label{lesson:RU-C409-polomka}



\begin{quote}
Before the new one: say the Russian for a driving licence, then the Russian for to fine.
\end{quote}

\subsection*{You'll want to know: \ru{поломка}}
\textbf{\ru{поломка}} (\emph{polómka}) — \textquotedblleft{}a breakdown, a fault\textquotedblright{}.

\begin{quote}
\ru{Из}-\ru{за} \ru{поломки} \ru{мы} \ru{приехали} \ru{домой} \ru{поздно}.
\end{quote}

\begin{quote}
\emph{Iz-za polómki my priyékhali domóy pózdno.}
\end{quote}

\begin{quote}
Because of the breakdown we got home late.
\end{quote}

Feminine. Built on \emph{\ru{поломать}} — to break.

\begin{cousinweb}
\emph{\ru{поломка}} is built on \emph{\ru{поломать}} 'to break', from \emph{\ru{ломать}} 'to break'; EtymDB lists \emph{\ru{ломать}} without a deeper source.
\end{cousinweb}

\begin{grammarlens}[title={What you've built}]
One more word for telling, explaining and describing at B1.
\end{grammarlens}

\subsection*{Your turn}
\begin{itemize}
\raggedright
  \item \emph{Say it:} \emph{polómka}
  \item \emph{Say it:} \emph{polómka}, once more
  \item \emph{Say it:} a sentence of your own with \emph{polómka}
  \item \emph{Recall:} say the Russian for a driving licence, then the Russian for to fine, then say \emph{polómka} again
\end{itemize}

\begin{tcolorbox}[breakable,colback=teal!4,colframe=teal!35!black,title={Before you move on}]
What does \ru{поломка} mean? (\textquotedblleft{}A breakdown, a fault\textquotedblright{}.) Say it once more.
\end{tcolorbox}
\section[\texorpdfstring{\ru{эвакуатор} (evakuátor)}{эвакуатор (evakuátor)}]{\ru{эвакуатор} (evakuátor) — a tow truck}

\label{lesson:RU-C409-evakuator}



\begin{quote}
Before the new one: say the Russian for to fine, then the Russian for a breakdown.
\end{quote}

\subsection*{You'll want to know: \ru{эвакуатор}}
\textbf{\ru{эвакуатор}} (\emph{evakuátor}) — \textquotedblleft{}a tow truck\textquotedblright{}.

\begin{quote}
\ru{Эвакуатор} \ru{приехал} \ru{через} \ru{час}.
\end{quote}

\begin{quote}
\emph{Evakuátor priyékhal chérez chas.}
\end{quote}

\begin{quote}
The tow truck came an hour later.
\end{quote}

Masculine. Built on \emph{\ru{эвакуировать}} — to evacuate, to remove.

\begin{cousinweb}
\emph{\ru{эвакуатор}} is built on \emph{\ru{эвакуировать}} 'to evacuate, to clear away', which EtymDB lists without a source; it is usually traced to Latin \emph{ēvacuāre} 'to empty out', on \emph{vacuus} 'empty' (Lewis \& Short).
\end{cousinweb}

\begin{grammarlens}[title={What you've built}]
One more word for telling, explaining and describing at B1.
\end{grammarlens}

\subsection*{Your turn}
\begin{itemize}
\raggedright
  \item \emph{Say it:} \emph{evakuátor}
  \item \emph{Say it:} \emph{evakuátor}, once more
  \item \emph{Say it:} a sentence of your own with \emph{evakuátor}
  \item \emph{Recall:} say the Russian for to fine, then the Russian for a breakdown, then say \emph{evakuátor} again
\end{itemize}

\begin{tcolorbox}[breakable,colback=teal!4,colframe=teal!35!black,title={Before you move on}]
What does \ru{эвакуатор} mean? (\textquotedblleft{}A tow truck\textquotedblright{}.) Say it once more.
\end{tcolorbox}
\section[\texorpdfstring{\ru{буксировать} (buksírovat)}{буксировать (buksírovat)}]{\ru{буксировать} (buksírovat) — to tow}

\label{lesson:RU-C409-buksirovat}



\begin{quote}
Before the new one: say the Russian for a breakdown, then the Russian for a tow truck.
\end{quote}

\subsection*{You'll want to know: \ru{буксировать}}
\textbf{\ru{буксировать}} (\emph{buksírovat}) — \textquotedblleft{}to tow\textquotedblright{}.

\begin{quote}
\ru{Сосед} \ru{буксировал} \ru{нашу} \ru{машину} \ru{до} \ru{дома}.
\end{quote}

\begin{quote}
\emph{Soséd buksíroval náshu mashínu do dóma.}
\end{quote}

\begin{quote}
The neighbour towed our car home.
\end{quote}

Imperfective; perfective \emph{\ru{отбуксировать}}. Present \emph{\ru{буксирую}}.

\begin{cousinweb}
\emph{\ru{буксировать}} is built on \emph{\ru{буксир}} 'a tug, a towrope'. \emph{\ru{Буксир}} is usually traced to a Dutch sailors' word for towing; EtymDB lists it without a source, so the line is hedged.
\end{cousinweb}

\begin{grammarlens}[title={What you've built}]
One more word for telling, explaining and describing at B1.
\end{grammarlens}

\subsection*{Your turn}
\begin{itemize}
\raggedright
  \item \emph{Say it:} \emph{buksírovat}
  \item \emph{Say it:} \emph{buksírovat}, once more
  \item \emph{Say it:} a sentence of your own with \emph{buksírovat}
  \item \emph{Recall:} say the Russian for a breakdown, then the Russian for a tow truck, then say \emph{buksírovat} again
\end{itemize}

\begin{tcolorbox}[breakable,colback=teal!4,colframe=teal!35!black,title={Before you move on}]
What does \ru{буксировать} mean? (\textquotedblleft{}To tow\textquotedblright{}.) Say it once more.
\end{tcolorbox}
\section[\texorpdfstring{\ru{домкрат} (domkrát)}{домкрат (domkrát)}]{\ru{домкрат} (domkrát) — a jack (for a car)}

\label{lesson:RU-C409-domkrat}



\begin{quote}
Before the new one: say the Russian for a tow truck, then the Russian for to tow.
\end{quote}

\subsection*{You'll want to know: \ru{домкрат}}
\textbf{\ru{домкрат}} (\emph{domkrát}) — \textquotedblleft{}a jack (for a car)\textquotedblright{}.

\begin{quote}
\ru{Без} \ru{домкрата} \ru{колесо} \ru{не} \ru{поменять}.
\end{quote}

\begin{quote}
\emph{Bez domkráta kolesó ne pomenyát.}
\end{quote}

\begin{quote}
You cannot change a wheel without a jack.
\end{quote}

Masculine; the stress is on the last syllable.

\begin{cousinweb}
\emph{\ru{домкрат}} was borrowed from Dutch \emph{dommekracht} (EtymDB), a word usually explained as \emph{dom} 'dumb' and \emph{kracht} 'strength': a tool that lifts by plain force.
\end{cousinweb}

\begin{grammarlens}[title={What you've built}]
One more word for telling, explaining and describing at B1.
\end{grammarlens}

\subsection*{Your turn}
\begin{itemize}
\raggedright
  \item \emph{Say it:} \emph{domkrát}
  \item \emph{Say it:} \emph{domkrát}, once more
  \item \emph{Say it:} a sentence of your own with \emph{domkrát}
  \item \emph{Recall:} say the Russian for a tow truck, then the Russian for to tow, then say \emph{domkrát} again
\end{itemize}

\begin{tcolorbox}[breakable,colback=teal!4,colframe=teal!35!black,title={Before you move on}]
What does \ru{домкрат} mean? (\textquotedblleft{}A jack (for a car)\textquotedblright{}.) Say it once more.
\end{tcolorbox}
\section[\texorpdfstring{\ru{автосервис} (avtosérvis)}{автосервис (avtosérvis)}]{\ru{автосервис} (avtosérvis) — a car repair shop, a garage}

\label{lesson:RU-C409-avtoservis}



\begin{quote}
Before the new one: say the Russian for to tow, then the Russian for a jack.
\end{quote}

\subsection*{You'll want to know: \ru{автосервис}}
\textbf{\ru{автосервис}} (\emph{avtosérvis}) — \textquotedblleft{}a car repair shop, a garage\textquotedblright{}.

\begin{quote}
\ru{Мы} \ru{оставили} \ru{машину} \ru{в} \ru{автосервисе} \ru{на} \ru{два} \ru{дня}.
\end{quote}

\begin{quote}
\emph{My ostávili mashínu v avtosérvise na dva dnya.}
\end{quote}

\begin{quote}
We left the car at the garage for two days.
\end{quote}

Masculine. In speech, often \emph{\ru{сервис}}.

\begin{cousinweb}
\emph{\ru{автосервис}} is built from \emph{\ru{авто}-}, short for \emph{\ru{автомобиль}} 'a car', and \emph{\ru{сервис}} 'service', which was borrowed from English \emph{service} (EtymDB). EtymDB traces the English word through Old French to Latin \emph{servitium}.
\end{cousinweb}

\begin{grammarlens}[title={What you've built}]
One more word for telling, explaining and describing at B1.
\end{grammarlens}

\subsection*{Your turn}
\begin{itemize}
\raggedright
  \item \emph{Say it:} \emph{avtosérvis}
  \item \emph{Say it:} \emph{avtosérvis}, once more
  \item \emph{Say it:} a sentence of your own with \emph{avtosérvis}
  \item \emph{Recall:} say the Russian for to tow, then the Russian for a jack, then say \emph{avtosérvis} again
\end{itemize}

\begin{tcolorbox}[breakable,colback=teal!4,colframe=teal!35!black,title={Before you move on}]
What does \ru{автосервис} mean? (\textquotedblleft{}A car repair shop, a garage\textquotedblright{}.) Say it once more.
\end{tcolorbox}
